\section{Limitations and conclusion}
\label{sec:limits}

\subsection{Limitations}

This report does not provide a completed Lean implementation of Boolean
reflection, trace coding, strategy reflection, or the open bridge.  It provides
full human arguments, exact object-theory endpoints, a no-placeholder
structural kernel candidate, and a dependency-ordered path to those proofs.
The absence of a Lean executable in the artifact runtime prevented even the
small candidate from being elaborated, so no new kernel-checked theorem is
claimed.

The framework does not prove the deep producer theorems identified in the
packet catalogue.  In particular, it proves no general Faltings, Siegel,
Mordell--Weil, effective Roth, Chebotarev, or unrestricted analytic convergence
result.  It does not select exact branches through singularities, derive
moduli from arbitrary classical convergence proofs, or infer an enumeration
cutoff from finite cardinality.

The fixed-bounded-formula compiler is a schema indexed by host syntax.  This is
intentional.  Internalizing a universal interpreter would be a different and
substantially harder arithmetization project, and a total universal function
for all primitive-recursive programs is not the interface needed by the
applications.

Finally, ordinary \(\Prf\) proof terms may still be large.  Trace compression,
shared Lean definitions, and an optional \(\KPrf\) backend address the main
sources of blow-up, but empirical compilation will determine which
subgraphs should be flattened or translated.

\subsection{Conclusions}

The current \repo\ tree already contains the logical and arithmetic machinery
needed for a principled general framework.  The missing step is to use that
machinery at the right granularity.

A mathematical proof in iPRA should be organized around finite queries and
local contracts.  Auxiliary certificates should remain existentially sealed
unless something certificate-dependent must escape.  Exact invariant
observations should replace arbitrary representative choices; approximate
clouds should replace unstable labels; a positive gap should be the explicit
gate to exact naming.  Variable-length work should be represented by a whole
trace or finite strategy and reflected by a quantifier-free invariant.  Missing
height bounds, moduli, covers, or constants should appear as formulas in the
final theorem, not as prose in a ``future work'' section.

This organization produces a small, auditable compiler whose primitive steps
are already checked: substitution, weakening, simultaneous cut, existential
elimination, and qf induction.  It also explains how the strongest ideas of the
existing packet frameworks fit together.  FQNF is the packed named backend;
CBIP is a bounded-interaction front end; \(\KPrf\) is a compact qf backend;
SCR is the ordinary logical spine.

The immediate practical recommendation is therefore narrow: elaborate the
supplied structural kernel, implement the search and trace adapters against
\code{IPRADerivations}, and migrate the four pilots before designing another
large proof-graph language.  Success at those gates would turn a substantial
class of current ``human-complete but uncompiled'' packets into literal
\(\Prf\;\PRAAx\) derivations, while leaving the genuinely deep mathematical
boundaries visible and honest.
